\section{Introduction}
Due to the quadratic time complexity of attention, its efficient implementation becomes crucial as sequence lengths increase in real-world applications~\cite{jiang2024minference,zhang2025sage}. Several strategies have been developed to mitigate the computational demands of attention —such as (1) \textit{Linear attention} methods~\cite{wang2020linformer,choromanski2020rethinking,yu2022metaformer,katharopoulos2020transformers} that reduce complexity to $O(N)$ and (2) \textit{Sparse attention} methods~\cite{liu2021swin,chu2021twins,liuniformer,xiao2023efficient,xiao2024infllm,chen2023longlora,jiang2024minference,venkataramanan2023skip,gao2024seerattention,moaattention,zhang2025spargeattn,xi2025sparse,yang2025sparse,zhang2025spargeattn_wksp,zhang2025fast,zhang2025faster} that selectively process parts of the context — these methods are only suitable for a limited range of models and tasks. The widely used attention methods optimize attention by exploiting hardware capacities, such as FlashAttention V1, V2, V3~\cite{dao2022flashattention,dao2023flashattention,shah2024flashattention}, xformers~\cite{xFormers2022}, and SageAttention~\cite{2024sageattention,zhang2025sageattention3,zhang2025sageattention2++,zhang2025sageattention2_wksp}. These works do not omit computations across the entire sequence and achieve impressive speed and accuracy performance across various tasks.

\noindent\textbf{{Motivation}}. For the two matrix multiplication (Matmul) operations in attention: $QK^\top$ and $\widetilde{P}V$, SageAttention accelerates them by quantizing the $QK^\top$ to INT8 and uses FP16 Matmul with FP16 accumulators for $\widetilde{P}V$. Moreover, to maintain attention accuracy, SageAttention proposes smoothing $K$ by eliminating its channel-wise outliers. SageAttention achieves 2 $\times$ speedup compared with FlashAttention2 and is the first quantized attention that incurs a negligible end-to-end metrics loss across language, image, and video generation models. However, SageAttention has two weaknesses. \textbf{(W1)} INT8 Matmul achieves only half the speed of INT4. \textbf{(W2)} FP16 Matmul with FP16 accumulators provides a speedup only on RTX 4090 and RTX 3090 GPUs.
To leverage the faster INT4 tensor cores for $QK^\top$ and using a method that can accelerate $\widetilde{P}V$ on a broader range of GPUs, we propose to quantize $Q, K$ to INT4 and $\widetilde{P}, V$ to FP8.

\noindent\textbf{{Challenges}}. Quantizing $Q, K$ to INT4 and $\widetilde{P}, V$ to FP8 presents significant challenges. For example, when only per-tensor quantizing $Q, K$ to INT4, the text-to-video model CogvideoX will generate a completely blurry video, and Llama3 only achieves a random-guessing-level accuracy of 25\% on MMLU. After an in-depth investigation, we identified three primary challenges: \textbf{(C1)} The numerical range of INT4 is quite limited compared to FP16 or INT8, leading to significant quantization errors when $Q$ and $K$ have some abnormal values. 
\textbf{(C2)} We discover that the FP32 accumulator designed for FP8 matrix multiplication in the tensor core (\texttt{mma.f32.f8.f8.f32}) is actually FP22, specifically with 1 sign bit, 8 exponent bits, and 13 mantissa bits. This will lead to an accuracy loss of $\widetilde{P}V$. 

\noindent\textbf{{Our approach}}. 
To address \textbf{(C1)}, we first discover that the per-block quantization of $Q, K$ in SageAttention does not achieve sufficient accuracy for INT4 quantization. Simultaneously, to avoid the extra latency caused by per-token dequantization, where each GPU thread corresponds to multiple quantization scales, we propose an adequately precise quantization method that incurs no additional dequantization overhead. Specifically, we introduce a per-thread quantization approach based on the mapping between the GPU threads and the memory layout of matrices as dictated by the PTX \texttt{mma} instructions. This method groups tokens corresponding to the same thread for quantization and dequantization, ensuring that each thread is associated with a single quantization scale. This approach achieves much better accuracy performance than per-block quantization with no additional dequantization overhead.
Second, for the significant channel-wise outliers in matrices $Q$ and $K$, we adopt smoothing $K$ in SageAttention and further propose an effective method to remove these outliers in $Q$. Specifically, we propose subtracting the average value of the channel dimension of $Q$, referred to as $\overrightarrow Q_m$. Then, we add $\overrightarrow Q_mK$ after the $QK^\top$ Matmul to ensure the correctness of the attention computation.
To address \textbf{(C2)}, the accuracy loss caused by the 22-bit accumulator for FP8 Matmul of $\widetilde{P}V$, we propose a two-level accumulation strategy that uses an FP32 buffer to accumulate the values from the 22-bit accumulator after each block Matmul of $\widetilde{P}V$. This method confines the errors to the block range.
Additionally, we design an optional technique to enhance the accuracy of the 22-bit accumulator. Specifically, we could smooth $V$ by subtracting the average value of its channel dimension and adding the subtracted item to the attention output to maintain the correctness.


\noindent\textbf{{Performance.}} We offer a high-performance implementation of \our on RTX4090 and L20 GPUs. This implementation achieves a peak performance of \textbf{481 TOPS} on the RTX4090, outperforming FlashAttention2 and xformers by approximately 3x and 4.5x, respectively.
To support NVIDIA Hopper GPUs, which lack native INT4 tensor core support, we additionally provide \texttt{SageAttention2-8b}\xspace, a variant that quantizes $Q,K$ to INT8. FlashAttention3, in contrast, is tailored to and only compatible with the Hopper architecture. Moreover, \texttt{SageAttention2-8b}\xspace matches the speed of FlashAttention3(fp8) on Hopper GPUs, while delivering much better accuracy. For example, on popular video generation models, our method does not compromise end-to-end accuracy, whereas FlashAttention3(fp8) brings noticeable degradation, as visualized in Fig.~\ref{fig:video_example1} and \ref{fig:video_example}.
We extensively evaluate the end-to-end metrics of state-of-the-art text, image, and video generation models. \our can accelerate models in a plug-and-play way with negligible loss in end-to-end metrics.


